\section{Proof of the uniform trace law}\label{app:general-trace}

The proof combines a full-fiber unitary transport, scalar Hilbert-square
corners and a Carleman--Mellin comparison. The classical comparison
methods are in Otte~\cite[Lemmas~3.1--3.3]{Otte2024},
Fedele--Gebert~\cite[proof of Corollary~1.5]{FedeleGebert2019} and
Laptev--Safarov~\cite[Theorem~1.2]{LaptevSafarov1996}; all estimates used
here are proved below.

We use Section~\ref{sec:crosswindow} and
Proposition~\ref{T4:prop:bloch}: dilation identifies
$\mathcal K_\tau$ with $P_{(-\tau,\tau)}\mathcal G P_{(-\tau,\tau)}$,
where $P_I$ denotes multiplication by $\mathbf1_I$ and
\begin{equation}\label{rev5:eq:G}
 G(X,Y)=\frac{\sin(2\pi X)-\sin(2\pi Y)}{\pi(X-Y)},\qquad
 \mathcal G=\mathcal G^*,\quad \|\mathcal G\|=1,\quad
 \mathcal G^2=\mathcal F^*\mathbf1_{(-1,1)}\mathcal F.
\end{equation}
The global Fourier transform $\mathcal F$ retains the positive exponent
of Section~\ref{sec:crosswindow}.  Bloch coefficients below use a negative
exponent, as in \eqref{T4:eq:symbol}; the angular Fourier transform used
after the Mellin change of variables will be defined separately.
All Schatten norms are denoted by $\|\cdot\|_p$; in particular,
$\mathfrak S_1$ is the trace class.  Operators compressed to a subspace
are extended by zero when they are compared on a common ambient space.

\subsection{Trace comparisons and bounded endpoint changes}
\label{rev5:subsec:trace-comparisons}

\begin{lemma}[Scalar trace perturbation]\label{rev5:lem:trace-lip}
Let $A,B$ be self-adjoint trace-class operators with spectra in a compact
interval containing zero.  If $f$ is real and Lipschitz on that interval,
with $f(0)=0$, then
\begin{equation}\label{rev5:eq:trace-lip}
 |\operatorname{tr}f(A)-\operatorname{tr}f(B)|
 \leq\operatorname{Lip}(f)\|A-B\|_1.
\end{equation}
\end{lemma}

\begin{proof}
For matrices, differentiate a polynomial along $A_t=A+t(B-A)$:
$\frac d{dt}\tr p(A_t)=\tr(p'(A_t)(B-A))$. Integration gives the
bound with $\|p'\|_\infty$. Uniform polynomial approximation of a
continuous derivative, then convolution of a Lipschitz extension,
gives every Lipschitz $f$ without increasing its Lipschitz constant.
For trace-class operators apply the matrix result on the sum of their
finite spectral subspaces and pass to the limit. Both operator tails
and scalar trace tails converge because
$|f(\lambda)|\le\operatorname{Lip}(f)|\lambda|$.
This is a scalar trace inequality, not an operator-Lipschitz assertion.
\end{proof}

\begin{lemma}[Fixed strips]\label{rev5:lem:strips}
For an interval $E$ of length $\ell$, put
\begin{equation}\label{rev5:eq:strip-function}
 b(\ell)=\sqrt{2\ell}\int_0^1e^{\pi\ell u^2}\,du.
\end{equation}
Then $\|P_E\mathcal G\|_1\leq b(\ell)
\leq\sqrt{2\ell}e^{\pi\ell}$.  If $I\mathbin\triangle J$ consists of
at most two intervals, each of length at most $h$, then
\begin{equation}\label{rev5:eq:strips}
 \|P_I\mathcal G P_I-P_J\mathcal G P_J\|_1\leq4b(h).
\end{equation}
In particular, writing $A_I=P_I\mathcal G P_I$, every interval
$I\subseteq\mathbb R$ and every $a_1,a_2\in\mathbb R$ satisfy
\begin{equation}\label{rev5:eq:translated-intervals}
 \|A_{I+a_1}-A_{I+a_2}\|_1\leq4b(|a_1-a_2|).
\end{equation}
\end{lemma}

\begin{proof}
The restricted Fourier operator $R_E:L^2(-1,1)\to L^2(E)$ with kernel
$e^{2\pi iX\xi}$ satisfies
$R_ER_E^*=P_E\mathcal G^2P_E
=(P_E\mathcal G)(P_E\mathcal G)^*$, by the symmetry of the frequency
interval in \eqref{rev5:eq:G}.  It has the same nonzero singular values
as $P_E\mathcal G$.  Translation of $E$ gives a unitary factor in the
frequency variable, so assume $E=[-\ell/2,\ell/2]$.  Expand its kernel as
$\sum_{j\geq0}(2\pi i)^jX^j\xi^j/j!$.  The $j$th rank-one term has
trace norm
\[
 \frac{(2\pi)^j}{j!}\|X^j\|_{L^2(E)}\|\xi^j\|_{L^2(-1,1)}
 =\sqrt{2\ell}\,\frac{(\pi\ell)^j}{j!(2j+1)}.
\]
Their sum is $b(\ell)$; uniform kernel convergence identifies the
trace-norm sum with $R_E$. For $D=P_I-P_J$, the identity
$P_I\mathcal G P_I-P_J\mathcal G P_J=D\mathcal G P_I+P_J\mathcal G D$
bounds the difference by $2\|D\mathcal G\|_1$. Decompose $D$ into its
at most two signed strip projections. Monotonicity of $b$ gives
\eqref{rev5:eq:strips}; two translates have strips of length at most
$|a_1-a_2|$, including when they are disjoint.
\end{proof}

\begin{corollary}[Odd tests and integer parameters]
\label{rev5:cor:odd-endpoints}
For every odd real Lipschitz function $f$ and every $\tau>0$,
\begin{equation}\label{rev5:eq:odd-bound}
 |\operatorname{tr}f(\mathcal K_\tau)|
 \leq2b(1/2)\operatorname{Lip}(f)
 \leq2e^{\pi/2}\operatorname{Lip}(f).
\end{equation}
For $\tau\geq2$ and a nearest integer $q$ to $\tau$,
\begin{equation}\label{rev5:eq:integer-transfer}
 |\operatorname{tr}f(\mathcal K_\tau)
       -\operatorname{tr}f(\mathcal K_q)|
 \leq4b(1/2)\operatorname{Lip}(f)
\end{equation}
for every real Lipschitz $f$ with $f(0)=0$.
\end{corollary}

\begin{proof}
The kernel obeys $G(X+1/2,Y+1/2)=-G(X,Y)$.  Thus the compression to
$I+1/2$ is unitarily equivalent to the negative of the compression to
$I$.  Their symmetric difference has at most two components, each of
length at most $1/2$.  Lemmas~\ref{rev5:lem:trace-lip}
and~\ref{rev5:lem:strips}, followed by division by two, prove
\eqref{rev5:eq:odd-bound}.  For \eqref{rev5:eq:integer-transfer}, the
two endpoints of $(-\tau,\tau)$ each move by at most $1/2$ when replaced
by those of $(-q,q)$.
\end{proof}

\subsection{A periodic gauge to a three-dimensional symbol}
\label{rev5:subsec:gauge}

On the Bloch fiber $\mathscr H=L^2(0,1)$ set
$e_j(t)=e^{2\pi ijt}$ and $W_\theta=M_{e^{2\pi i\theta t}}$.
Writing $X_{j,k}=|e_j\rangle\langle e_k|+|e_k\rangle\langle e_j|$,
the symbol of Proposition~\ref{T4:prop:bloch} is
\[
 \sigma(\theta):=\varsigma(\theta)=W_\theta^*X_{1,0}W_\theta,
 \qquad 0\leq\theta\leq1.
\]
Define
\begin{equation}\label{rev5:eq:star}
 c_\theta=\cos(\pi\theta/2),\quad s_\theta=\sin(\pi\theta/2),\quad
 w_\theta=c_\theta e_1+s_\theta e_{-1},\qquad
 a(\theta)=|e_0\rangle\langle w_\theta|
             +|w_\theta\rangle\langle e_0|.
\end{equation}
Both symbols have eigenvalues $1,-1,0$, with simple nonzero eigenvalues,
and the same endpoint operators:
$a(0)=\sigma(0)=X_{1,0}$ and $a(1)=\sigma(1)=X_{0,-1}$.
The range of $a$ lies in the fixed space
$E=\operatorname{span}\{e_{-1},e_0,e_1\}$.

\begin{lemma}[Full-fiber gauge]\label{rev5:lem:gauge}
There is a unitary path $V_\theta$ on $\mathscr H$ such that
\begin{equation}\label{rev5:eq:gauge}
 V_0=V_1=I,\qquad
 V-I\in C^\infty([0,1];\mathfrak S_1),\qquad
 V_\theta^*\sigma(\theta)V_\theta=a(\theta).
\end{equation}
Its periodic extension is continuous; matching derivatives at the two
ends is not required.
\end{lemma}

\begin{proof}
For either $b=\sigma$ or $b=a$, form its three spectral projections
\[
 P_+=(b^2+b)/2,\qquad P_-=(b^2-b)/2,\qquad P_0=I-b^2.
\]
The rank-one projections $P_\pm$ are smooth in trace norm, by their
Riesz formulas and the rank-two bound for their differences; derivatives
of $P_0$ are minus their sums. Thus
$K_b=\sum_{j\in\{+,-,0\}}P_j'P_j$ is smooth in trace norm and
skew-adjoint. The projection identities give $[P_j,K_b]=-P_j'$,
so $U_b'=K_bU_b$, $U_b(0)=I$, has a unitary solution satisfying
\[
 U_b(\theta)^*P_j(\theta)U_b(\theta)=P_j(0),\qquad
 U_b(\theta)-I=\int_0^\theta K_b(t)U_b(t)\,dt.
\]
The Bochner integral and differential equation make $U_b-I$ smooth
in trace norm. Put $S=U_\sigma$, $T=U_a$, $M=T_1^*S_1$. Endpoint
agreement makes $M$ commute with $X_{1,0}$, and $M-I$ is trace class.
Compact-normal spectral calculus gives a commuting skew-adjoint
logarithm $L$: choose eigenangles $\phi\in(-\pi,\pi]$, with zero on
the eigenspace of one, and use
$|\phi|\le(\pi/2)|e^{i\phi}-1|$ for trace summability.
Then $V_\theta=S_\theta e^{-\theta L}T_\theta^*$ conjugates
$a(\theta)=T_\theta X_{1,0}T_\theta^*$ to $\sigma(\theta)$ and has
$V_0=V_1=I$. The construction includes the infinite zero eigenspace.
\end{proof}

For a bounded operator-valued symbol $b$, use the conventions
\begin{equation}\label{rev5:eq:laurent}
 \widehat b(k)=\int_0^1b(\theta)e^{-2\pi ik\theta}\,d\theta,
 \qquad (\mathcal L(b)x)_j=\sum_l\widehat b(j-l)x_l,
 \qquad T_n(b)=P_n\mathcal L(b)P_n,
\end{equation}
where $P_n$ selects cells $0,\ldots,n-1$ and acts as identity on each
fiber.  Laurent multiplication is multiplicative and preserves adjoints.

\begin{lemma}[A continuous seam]\label{rev5:lem:seam}
Suppose $F\in C^3([0,1];\mathfrak S_1)$ and $F(0)=F(1)$.  Define
\begin{equation}\label{rev5:eq:seam-constant}
 \mathfrak c(F)=\frac{\|F'(1)-F'(0)\|_1}{8}
 +\frac{\|F''(1)-F''(0)\|_1+
             \int_0^1\|F'''(\theta)\|_1\,d\theta}{12\pi}.
\end{equation}
Then $\|[P_n,\mathcal L(F)]\|_1\leq\mathfrak c(F)$ for every $n$.
The same assertion applies to scalar symbols, using absolute values
in \eqref{rev5:eq:seam-constant}.
\end{lemma}

\begin{proof}
For $k\ne0$, three integrations by parts give
\[
 \widehat F(k)=\frac{D}{4\pi^2k^2}+R_k,
 \quad D=F'(1)-F'(0),\quad
 R_k=\frac{-D_2+\int_0^1F'''(\theta)e^{-2\pi ik\theta}\,d\theta}
            {(2\pi ik)^3},
\]
where $D_2=F''(1)-F''(0)$.  Thus
\[
 \sum_{k\ne0}|k|\|R_k\|_1
 \leq\frac{\|D_2\|_1+\int_0^1\|F'''(\theta)\|_1\,d\theta}{24\pi}.
\]
For the principal term use the positive Hankel matrix
\[
 H^{(2)}_{jl}=(j+l+1)^{-2}
 =\int_0^1t^{j+l}(-\log t)\,dt,
 \qquad \|H^{(2)}\|_1=\sum_{j\geq0}(2j+1)^{-2}=\pi^2/8.
\]
The rank-one integral proves positivity and finite trace. Each end
block of the Laurent sequence $k^{-2}$ compresses $H^{(2)}$, with
row reversal at the right end. The commutator therefore has norm
at most $4\|H^{(2)}\|_1$, giving $\|D\|_1/8$ after tensoring.
For the remainder use the bilateral shift $\mathsf S$ and
$\|[P_n,\mathsf S^k\otimes R_k]\|_1=2\min(n,|k|)\|R_k\|_1$.
Summation proves \eqref{rev5:eq:seam-constant}; the zeroth coefficient
commutes with $P_n$ and the Fourier series converges in operator norm.
\end{proof}

Apply the lemma to $F=V-I$ and put $C_V=\mathfrak c(V-I)$.
The Laurent unitary $\mathscr V=\mathcal L(V)$ satisfies
$\mathcal L(a)=\mathscr V^*\mathcal L(\sigma)\mathscr V$.  The exact
identity, with $B=\mathcal L(\sigma)$,
\[
 \mathscr V^*P_nBP_n\mathscr V-P_n\mathscr V^*B\mathscr VP_n
 =(\mathscr V^*P_n-P_n\mathscr V^*)BP_n\mathscr V
   +P_n\mathscr V^*B(P_n\mathscr V-\mathscr VP_n)
\]
therefore has trace norm at most $2C_V$, since $\|B\|=1$.
Both compressions are trace class: for $\sigma$ this follows from the
original finite-interval operator and Proposition~\ref{prop:traces};
for $a$ the compression acts on $E^n$ and vanishes on its complement.
Lemma~\ref{rev5:lem:trace-lip} gives
\begin{equation}\label{rev5:eq:star-transfer}
 |\operatorname{tr}f(T_n(\sigma))
       -\operatorname{tr}_{E^n}f(T_n(a))|
 \leq2C_V\operatorname{Lip}(f).
\end{equation}
The gauge compressed to $P_n$ was not assumed unitary.

\begin{lemma}[A numerical bound for the full-fiber gauge]
\label{pubpair:lem:numerical-gauge}
For the path of Lemma~\ref{rev5:lem:gauge}, with the specified
eigenangle logarithm, the constant in
\eqref{rev5:eq:star-transfer} satisfies $C_V<65$.
\end{lemma}
\begin{proof}
On the whole fiber set
\[
 \mathsf D=M_{2\pi i(t-1/2)},\qquad
 P=|e_0\rangle\langle e_0|+|e_1\rangle\langle e_1|,\qquad
 Q=I-P,\qquad R=Q\mathsf DQ.
\]
Both $\mathsf D$ and $R$ are skew-adjoint, with norms at most $\pi$.
The scalar phase introduced by centering does not change the
spectral projections of $\sigma$, so
$P_j(\theta)=e^{-\theta\mathsf D}P_j(0)e^{\theta\mathsf D}$ and
\[
 K_\sigma(\theta)=e^{-\theta\mathsf D}K_0e^{\theta\mathsf D},
 \qquad
 K_0=\sum_jP_j(0)\mathsf DP_j(0)-\mathsf D.
\]
Indeed $P_j'=P_j\mathsf D-\mathsf DP_j$.  In the basis $(e_0,e_1)$,
$P\mathsf DP=\left(\begin{smallmatrix}0&1\\-1&0\end{smallmatrix}\right)$;
its expectations on $(e_0\pm e_1)/\sqrt2$ vanish.  Thus
\[
 K_0=Q\mathsf DQ-\mathsf D
    =-P\mathsf DP-P\mathsf DQ-Q\mathsf DP.
\]
Its range lies in $P\mathscr H+\operatorname{ran}(Q\mathsf DP)$, of
dimension at most four.  The three displayed blocks are orthogonal
in Hilbert--Schmidt inner product.  Direct integration gives
$\|P\mathsf DP\|_2^2=2$ and
$\operatorname{tr}(P\mathsf D^*\mathsf DP)=2\pi^2/3$, whence
\[
 \|K_0\|_2^2=4\pi^2/3-2,\qquad
 k_0:=\|K_0\|_1\le2\sqrt{4\pi^2/3-2}<67/10.
\]
The original transport is exactly
$S_\theta=e^{-\theta\mathsf D}e^{\theta R}$: its derivative is
$e^{-\theta\mathsf D}K_0e^{\theta R}=K_\sigma(\theta)S_\theta$,
and $S_0=I$.  Differentiating this formula further yields
\[
 \|S_\theta^{(j)}\|_1\le(2\pi)^{j-1}k_0,\qquad j=1,2,3.
\]
Only this differentiated product is estimated in trace norm; its
two exponential factors need not separately be trace-class
perturbations of the identity.

For the star symbol put
$J=(\pi/2)(|e_{-1}\rangle\langle e_1|
          -|e_1\rangle\langle e_{-1}|)$.
The unitary $e^{\theta J}$ fixes $e_0$ and rotates $e_1$ to $w_\theta$.
All three diagonal spectral blocks $P_j(\theta)JP_j(\theta)$
vanish, including the zero block, so $K_a=J$ and
$T_\theta=e^{\theta J}$.  Since $J$ has rank two,
\[
 \|(T_\theta^*)^{(j)}\|\le(\pi/2)^j,\qquad
 \|(T_\theta^*)^{(j)}\|_1=\pi(\pi/2)^{j-1},\qquad
 \|T_1-I\|_1=2\sqrt2.
\]
For $M=T_1^*S_1$ and the logarithm $L$ used in
Lemma~\ref{rev5:lem:gauge}, the eigenangle inequality in that proof
gives
\[
 \|L\|\le\pi,\qquad
 \ell_0:=\|L\|_1
 \le\frac{\pi}{2}\|M-I\|_1
 \le\frac{\pi}{2}(k_0+2\sqrt2)<15.
\]
Here $\|S_1-I\|_1\le k_0$ follows by integration.  No finite-rank
assumption on $M-I$ or $L$ is used.  The numerical bounds follow
already from $\pi<22/7$ and $\sqrt2<99/70$.

Apply Leibniz to $V_\theta=S_\theta e^{-\theta L}T_\theta^*$.
In each term use the trace norm of the first differentiated factor
and operator norms for the other factors.  The resulting uniform
bounds on the first three derivatives are
\begin{equation}\label{pubpair:eq:gauge-derivatives}
 \begin{aligned}
 v_1&=k_0+\ell_0+\pi,\\
 v_2&=5\pi k_0+2\pi\ell_0+\pi^2/2,\\
 v_3&=(79/4)\pi^2k_0+(13/4)\pi^2\ell_0+\pi^3/4.
 \end{aligned}
\end{equation}
Explicitly, for $j=1,2,3$ the product sum is bounded by
\[
 \frac{k_0}{2\pi}\big[(7\pi/2)^j-(3\pi/2)^j\big]
 +\frac{\ell_0}{\pi}\big[(3\pi/2)^j-(\pi/2)^j\big]
 +\pi(\pi/2)^{j-1},
\]
which gives \eqref{pubpair:eq:gauge-derivatives}.
The two endpoint differences in
\eqref{rev5:eq:seam-constant} are bounded by $2v_1,2v_2$,
and the third-derivative integral by $v_3$.  Consequently
\[
 \begin{split}
 C_V&\le\frac{v_1}{4}+\frac{2v_2+v_3}{12\pi}\\
 &\le k_0\left(\frac{13}{12}+\frac{79\pi}{48}\right)
    +\ell_0\left(\frac7{12}+\frac{13\pi}{48}\right)
    +\frac{\pi}{3}+\frac{\pi^2}{48}\\
 &<\frac{760709}{11760}<65.
 \end{split}
\]
\end{proof}

\subsection{The scalar defect and the test function}
\label{rev5:subsec:scalar-defect}

Grouping the center and leaves of \eqref{rev5:eq:star} gives
\[
 T_n(a)\simeq\begin{pmatrix}0&B_n\\B_n^*&0\end{pmatrix},
 \qquad B_n=(T_n(s),T_n(c)),\qquad
 H_n=B_nB_n^*=T_n(c)^2+T_n(s)^2.
\]
Restricted to $E^n$ its eigenvalues are $\pm s_j(B_n)$, $1\leq j\leq n$, and $n$
additional zeros.  For a real scalar symbol $w$ write
\[
 D_n(w)=P_n\mathcal L(w)(I-P_n)\mathcal L(w)P_n
        =T_n(w^2)-T_n(w)^2\geq0.
\]
The functions $u=(c+s)/\sqrt2$, $v=(c-s)/\sqrt2$ satisfy
\begin{equation}\label{rev5:eq:scalar-defect}
 H_n=I-D_n(u)-D_n(v).
\end{equation}
Indeed $c^2+s^2=1$ gives $H_n=I-D_n(c)-D_n(s)$, and the
parallelogram law for $A_w=P_n\mathcal L(w)(I-P_n)$ gives
$D_n(u)+D_n(v)=D_n(c)+D_n(s)$. This rotation leaves only $v$
discontinuous at the periodic seam. Since $T_n(a)$ is a contraction,
$0\le H_n\le I$.

Set $v_0(\theta)=(1-2\theta)/\sqrt2$ and $z=v-v_0$.  The endpoint
values of $u$ agree; those of $z$ are both zero.  These functions are
smooth on $[0,1]$, so Lemma~\ref{rev5:lem:seam} applies.  In particular,
$\|P_n\mathcal L(w)(I-P_n)\|_1\leq\mathfrak c(w)$ for $w=u,z$.
Expanding $D_n(v)-D_n(v_0)$ and using boundedness of the other factors
gives
\begin{equation}\label{rev5:eq:star-defect-error}
 \|D_n(u)+D_n(v)-D_n(v_0)\|_1\leq C_*,\qquad
 C_*:=\|u\|_\infty\mathfrak c(u)
       +(\|v\|_\infty+\|v_0\|_\infty)\mathfrak c(z).
\end{equation}
The constant is independent of $n$, $f$ and $\tau$.
For example, if $A_w=P_n\mathcal L(w)(I-P_n)$, then
$D_n(v)-D_n(v_0)=A_zA_v^*+A_{v_0}A_z^*$, which proves the stated bound.
Also
\begin{equation}\label{rev5:eq:half-bound}
 0\leq D_n(v_0)\leq T_n(v_0^2)\leq\tfrac12 I.
\end{equation}

\begin{lemma}[Paired test function]\label{rev5:lem:paired-test}
Let $f\in C^2([-1,1])$ be real, with $f(0)=0$, and put
\begin{equation}\label{rev5:eq:h}
 h(x)=f(\sqrt{1-x})+f(-\sqrt{1-x})-f(1)-f(-1),\qquad 0\leq x\leq1.
\end{equation}
Then $h(0)=0$, $\operatorname{Lip}_{[0,1]}(h)\leq\|f''\|_\infty$,
and $h\in C^2([0,1/2])$ with
$\|h''\|_{[0,1/2]}\leq2\|f''\|_\infty$.
Moreover
\begin{equation}\label{rev5:eq:paired-trace}
 \left|\operatorname{tr}f(T_n(a))-n[f(1)+f(-1)]
                    -\operatorname{tr}h(D_n(v_0))\right|
 \leq C_*\|f''\|_\infty.
\end{equation}
\end{lemma}

\begin{proof}
For $r=\sqrt{1-x}>0$,
\[
 h'(x)=\frac{f'(-r)-f'(r)}{2r},\qquad
 h''(x)=\frac{f''(r)+f''(-r)}{4r^2}
          -\frac{f'(r)-f'(-r)}{4r^3}.
\]
Integrating $f''$ from $-r$ to $r$ proves the first bound, including
the limit at $r=0$.  The second expression is bounded by
$\|f''\|_\infty/r^2$, hence by $2\|f''\|_\infty$ when $x\leq1/2$.
Signed singular-value pairing gives exactly
\[
 \operatorname{tr}f(T_n(a))
 =n[f(1)+f(-1)]+\operatorname{tr}h(D_n(u)+D_n(v)).
\]
Possible zero singular values cause no extra term because $f(0)=0$.
Apply Lemma~\ref{rev5:lem:trace-lip} and
\eqref{rev5:eq:star-defect-error}.  Only the global Lipschitz bound
is used in this first comparison.  All later model spectra lie in
$[0,1/2]$, where the second derivative is bounded.
\end{proof}

\subsection{The Hilbert and Carleman operators}
\label{rev5:subsec:hilbert-carleman}

We establish boundedness before using the square of the infinite Hilbert
matrix.  Let $\mathcal C$ denote the Carleman operator with kernel
$(u+v)^{-1}$ on $L^2(0,\infty)$.  The unitary map
\[
 (\mathscr Mf)(s)=e^{s/2}f(e^s)
\]
turns it into convolution with $[2\cosh(s/2)]^{-1}$.  For the angular
Fourier convention $\widehat k(\omega)=\int_{\mathbb R}k(s)e^{-i\omega s}\,ds$,
Euler's integral at $z=1/2-i\omega$ gives
\begin{equation}\label{rev5:eq:carleman-multiplier}
 \int_{\mathbb R}\frac{e^{-i\omega s}}{2\cosh(s/2)}\,ds
 =\int_0^\infty\frac{t^{-1/2-i\omega}}{1+t}\,dt
 =\pi\operatorname{sech}(\pi\omega).
\end{equation}
Thus $\mathcal C$ is bounded, positive, and has norm $\pi$.
The operator $\mathcal C_{\rm sh}$ with kernel $(x+y+1)^{-1}$ is
unitarily its compression to $u\geq1/2$, by translation.

Embed $\ell^2(\mathbb N_0)$ as unit-cell constants in $L^2(0,\infty)$.
The Hilbert matrix $\mathsf H_{jl}=(j+l+1)^{-1}$, initially on finitely
supported sequences, has embedded kernel
$(\lfloor x\rfloor+\lfloor y\rfloor+1)^{-1}$ and is zero on the
orthogonal complement.  Write $\widetilde{\mathsf H}$ for this embedding.

\begin{lemma}[A nuclear comparison]\label{rev5:lem:hilbert-carleman}
The Hilbert matrix extends to a bounded operator, and
\begin{equation}\label{rev5:eq:hilbert-carleman}
 \|\widetilde{\mathsf H}-\mathcal C_{\rm sh}\|_1\leq\Delta,
 \qquad \Delta:=\frac{2}{1-e^{-2}}.
\end{equation}
\end{lemma}

\begin{proof}
Put $r_\lambda(x)=e^{-\lambda\lfloor x\rfloor}$,
$s_\lambda(x)=e^{-\lambda x}$.  The difference of the two kernels is
represented by
\[
 \int_0^\infty e^{-\lambda}
   (|r_\lambda\rangle\langle r_\lambda|
     -|s_\lambda\rangle\langle s_\lambda|)\,d\lambda.
\]
Here
$\|r_\lambda\|^2=(1-e^{-2\lambda})^{-1}$,
$\|s_\lambda\|^2=(2\lambda)^{-1}\leq\|r_\lambda\|^2$, and
$\|r_\lambda-s_\lambda\|\leq\min(\lambda,1)\|r_\lambda\|$.  The integrand's
trace norm is therefore at most
$2e^{-\lambda}\min(\lambda,1)/(1-e^{-2\lambda})$.  For $0<\lambda\leq1$,
concavity gives $1-e^{-2\lambda}\geq\lambda(1-e^{-2})$; for $\lambda\geq1$
use $1-e^{-2\lambda}\geq1-e^{-2}$.  Integration bounds the nuclear norm by
$\Delta$.  Testing against compactly supported functions identifies
the kernel of the nuclear integral with the displayed difference.
Since $\mathcal C_{\rm sh}$ is bounded, this also proves boundedness
of the Hilbert matrix.  Only the difference of the two Laplace
integrals has been integrated in trace norm.
\end{proof}

The Fourier coefficients of $v_0$ are
$\widehat v_0(k)=(\sqrt2\pi i k)^{-1}$ for $k\ne0$.
Multiplication of the left and right exterior blocks in $D_n(v_0)$ now
gives the exact identity
\begin{equation}\label{rev5:eq:hilbert-corners}
 D_n(v_0)=\alpha(P_n\mathsf H^2P_n+J_nP_n\mathsf H^2P_nJ_n),
 \qquad \alpha=\frac1{2\pi^2},
\end{equation}
where $J_n$ reverses the $n$ coordinates.  Indeed, each exterior index
produces the denominator $j+l+1$, and the opposite Fourier signs
multiply to the positive coefficient $\alpha$.
By positivity of both terms and \eqref{rev5:eq:half-bound},
$\alpha\|\mathsf H P_nx\|^2\leq\|P_nx\|^2/2$.
Testing finitely supported vectors proves $\|\mathsf H\|\leq\pi$.
Consequently \eqref{rev5:eq:hilbert-carleman} also gives
\begin{equation}\label{rev5:eq:square-comparison}
 \|\widetilde{\mathsf H}^2-\mathcal C_{\rm sh}^2\|_1
 \leq2\pi\Delta.
\end{equation}
In particular all compressions of $\alpha\mathsf H^2$,
$\alpha\mathcal C_{\rm sh}^2$ and $\alpha\mathcal C^2$ have
spectrum in $[0,1/2]$.

\subsection{Separating the two endpoints}
\label{rev5:subsec:corners}

\begin{lemma}[Smooth compression]\label{rev5:lem:compression}
Let $X$ be bounded and positive, $P$ an orthogonal projection, $PXP$
trace class, and $(I-P)XP$ Hilbert--Schmidt.  If $h$ is real and
$C^2$ on $[0,\|X\|]$, with $h(0)=0$, then
\begin{equation}\label{rev5:eq:compression}
 |\operatorname{tr}h(PXP)-\operatorname{tr}Ph(X)P|
 \leq\tfrac12\|h''\|_\infty\|(I-P)XP\|_2^2.
\end{equation}
The first trace is taken on $P\mathscr H$.  This is the positive case
of the compression estimate in \cite[Theorem~1.2]{LaptevSafarov1996}.
\end{lemma}

\begin{proof}
Since $|h(X)|\leq\operatorname{Lip}(h)X$, the compressed traces exist.
Choose an eigenbasis $v_j$ of the compact compression $PXP$, including
its kernel, with eigenvalues $\lambda_j$.  Scalar Taylor expansion in
the spectral measure of $X$ gives
\[
 |\langle v_j,h(X)v_j\rangle-h(\lambda_j)|
 \leq\tfrac12\|h''\|_\infty
                 \langle v_j,(X-\lambda_j)^2v_j\rangle
 =\tfrac12\|h''\|_\infty\|(I-P)Xv_j\|^2.
\]
The linear term vanishes because
$\langle v_j,Xv_j\rangle=\lambda_j$.  Summing proves the estimate.
\end{proof}

Take $n=2m$, and split its cells into two equal halves with projections
$P,Q$.  With $A=\alpha P_n\mathsf H^2P_n$ and
$A'=J_nAJ_n$, the trace of $A$ on the second half and that of $A'$ on the
first half obey
\begin{equation}\label{rev5:eq:wrong-half}
 \operatorname{tr}QAQ
 =\alpha\sum_{j=m}^{2m-1}\sum_{l\geq0}(j+l+1)^{-2}\leq2\alpha,
 \qquad \operatorname{tr}PA'P\leq2\alpha.
\end{equation}
The inner sum is at most $(j+1)^{-1}+(j+1)^{-2}$, and its sum
over the $m$ indices is at most $m/(m+1)+m/(m+1)^2<1$.
Since $\|A\|,\|A'\|\le1/2$, positivity gives
\[
 \|PAQ\|_2^2\leq\|A\|\operatorname{tr}QAQ\leq\alpha,
 \qquad \|PA'Q\|_2^2\leq\alpha.
\]
Thus $\|P(A+A')Q\|_2^2\leq4\alpha$.
Apply Lemma~\ref{rev5:lem:compression} to both complementary blocks
of $A+A'=D_n(v_0)$; their errors sum to at most
$4\alpha\|h''\|_{[0,1/2]}$.  Remove the other corner in each diagonal
half using Lemma~\ref{rev5:lem:trace-lip} and
\eqref{rev5:eq:wrong-half}.  The remaining two corners are identical.
We obtain the quantitative estimate
\begin{equation}\label{rev5:eq:decouple}
 \left|\operatorname{tr}h(D_{2m}(v_0))
       -2\operatorname{tr}h(\alpha P_m\mathsf H^2P_m)\right|
 \leq4\alpha\bigl(\|h'\|_{[0,1/2]}+\|h''\|_{[0,1/2]}\bigr).
\end{equation}

Under the step-function embedding, $P_m$ is restriction to $(0,m)$.
Its added zero eigenvalues do not affect the trace.  By
\eqref{rev5:eq:square-comparison} and Lemma~\ref{rev5:lem:trace-lip},
replacing its Hilbert-square corner by the shifted Carleman-square
corner costs at most $2\pi\alpha\Delta\operatorname{Lip}(h)$.
After translation the latter interval is $I=[a,m+a]$, $a=1/2$.
The difference between the compression of the full square and the
square of the half-line compression is the positive operator
\[
 P_I\mathcal C^2P_I
 -P_I\mathcal C P_{[a,\infty)}\mathcal C P_I
 =P_I\mathcal C P_{(0,a)}\mathcal C P_I,
\]
whose trace is
\begin{equation}\label{rev5:eq:missing-interval}
 \int_a^{m+a}\int_0^a\frac{dv\,du}{(u+v)^2}
 =\log\frac{2(m+a)}{m+2a}\leq\log2.
\end{equation}
Thus, for one corner,
\begin{equation}\label{rev5:eq:corner-transfer}
 \left|\operatorname{tr}h(\alpha P_m\mathsf H^2P_m)
       -\operatorname{tr}h(\alpha P_I\mathcal C^2P_I)\right|
 \leq\alpha(2\pi\Delta+\log2)\operatorname{Lip}(h).
\end{equation}

\subsection{The Mellin trace and completion of the proof}
\label{rev5:subsec:mellin}

By \eqref{rev5:eq:carleman-multiplier}, the Mellin transform of
$\alpha\mathcal C^2$ is convolution on $L^2(\mathbb R)$ with multiplier
and kernel
\begin{equation}\label{rev5:eq:mellin-model}
 d(\omega)=\tfrac12\operatorname{sech}^2(\pi\omega),\qquad
 \ell(s)=\frac{s}{4\pi^2\sinh(s/2)},\qquad \ell(0)=\frac1{2\pi^2}.
\end{equation}
Here inverse angular Fourier transformation has the factor $1/(2\pi)$.
Write $D$ for this convolution operator.  The interval $I$ becomes an
interval of length $L=\log(2m+1)$.  If $R_L$ restricts to such an
interval, then
\begin{equation}\label{rev5:eq:mellin-boundary}
 \|(I-R_L)DR_L\|_2^2
 =\int_{\mathbb R}\min(L,|s|)|\ell(s)|^2\,ds
 \leq M_D:=\frac{3\zeta(3)}{\pi^4}.
\end{equation}
To evaluate $M_D$, expand $\sinh^{-2}(s/2)=4\sum_{j\geq1}je^{-js}$ for
$s>0$ and integrate $s^3$ term by term; the first moment
$\int|s|\,|\ell(s)|^2\,ds$ converges at both zero and infinity.

The positive compression $R_LDR_L$ is trace class, of trace
$L\ell(0)$.  Since $h(0)=0$ and
$|h(d(\omega))|\leq\operatorname{Lip}(h)d(\omega)$, the integral
of $|h(d)|$ is finite.  Fourier factorization of its positive and
negative parts gives
\[
 \operatorname{tr}R_Lh(D)R_L
 =\frac{L}{2\pi}\int_{\mathbb R}h(d(\omega))\,d\omega.
\]
Lemma~\ref{rev5:lem:compression} and
\eqref{rev5:eq:mellin-boundary} consequently yield
\begin{equation}\label{rev5:eq:mellin-trace}
 \left|\operatorname{tr}h(R_LDR_L)
 -\frac{L}{2\pi}\int_{\mathbb R}h(d(\omega))\,d\omega\right|
 \leq\tfrac12 M_D\|h''\|_{[0,1/2]}.
\end{equation}

Combining \eqref{rev5:eq:decouple}, \eqref{rev5:eq:corner-transfer}
and \eqref{rev5:eq:mellin-trace} gives two copies of the Mellin bulk
term.  Since $\int_{\mathbb R}d(\omega)\,d\omega=1/\pi$, replacing
$\log(2m+1)$ by $\log n$, $n=2m$, costs at most
$\operatorname{Lip}(h)/\pi^2$.  In particular, with
\begin{equation}\label{rev5:eq:model-constants}
 B_1=4\alpha+4\pi\alpha\Delta+2\alpha\log2+\pi^{-2},
 \qquad B_2=4\alpha+M_D,
\end{equation}
we have
\begin{equation}\label{rev5:eq:model-trace}
 \left|\operatorname{tr}h(D_n(v_0))
    -\frac{\log n}{\pi}\int_{\mathbb R}h(d(\omega))\,d\omega\right|
 \leq B_1\operatorname{Lip}_{[0,1]}(h)
       +B_2\|h''\|_{[0,1/2]}
 \quad(n\geq2\text{ even}).
\end{equation}
For the paired test \eqref{rev5:eq:h}, substitute
\[
 t=\frac{1+\tanh(\pi\omega)}2,
 \qquad d(\omega)=2t(1-t),\qquad
 d\omega=\frac{dt}{2\pi t(1-t)}.
\]
The coefficient in \eqref{rev5:eq:model-trace} is exactly
\begin{equation}\label{rev5:eq:coefficient}
 \frac1\pi\int_{\mathbb R}h(d(\omega))\,d\omega
 =\frac1{2\pi^2}\int_0^1
 \frac{f(r_t)+f(-r_t)-f(1)-f(-1)}{t(1-t)}\,dt
 =\varkappa(f).
\end{equation}
Here $r_t=\sqrt{1-2t(1-t)}=\sqrt{t^2+(1-t)^2}$ agrees with the $r_t$ of
\eqref{T10:eq:kappa}.  The integral converges; for instance,
$1-r_t=2t(1-t)/(1+r_t)$ gives
\begin{equation}\label{rev5:eq:kappa-bound}
 |\varkappa(f)|\leq\frac{2}{\pi^2}\|f'\|_\infty.
\end{equation}

We now finish Theorem~\ref{T10:conj:law}.  At an integer parameter
$q\geq2$, the interval $(-q,q)$ has $n=2q$ complete cells; integer
translation does not change the compression's unitary class.  Equations
\eqref{rev5:eq:star-transfer}, \eqref{rev5:eq:paired-trace},
\eqref{rev5:eq:model-trace}, and Lemma~\ref{rev5:lem:paired-test} prove
\begin{equation}\label{rev5:eq:whole-cell-law}
 \begin{split}
 &\left|\operatorname{tr}f(\mathcal K_q)
       -n[f(1)+f(-1)]-\varkappa(f)\log n\right|\\
 &\hspace{1cm}\leq2C_V\|f'\|_\infty
             +(C_*+B_1+2B_2)\|f''\|_\infty.
 \end{split}
\end{equation}
For real $\tau\geq2$ choose a nearest integer $q$ and apply
\eqref{rev5:eq:integer-transfer}.  The bulk changes by at most
$2\|f'\|_\infty$, since $|2q-2\tau|\leq1$ and $f(0)=0$.
Also $1<2q/\tau\leq5/2$, so replacing $\log(2q)$ by $\log\tau$
costs at most $2\pi^{-2}\log(5/2)\|f'\|_\infty$, by
\eqref{rev5:eq:kappa-bound}.  Thus the estimate holds for every constant satisfying
\begin{equation}\label{rev5:eq:final-constant}
 C_{\mathrm{tr}}\ge2C_V+4b(1/2)+2+\frac{2\log(5/2)}{\pi^2}
       +C_*+B_1+2B_2.
\end{equation}
All quantities on the right are fixed independently of $f$ and
$\tau$, proving the uniform $C^2$ bound. No optimal constant is claimed.

\begin{remark}[An admissible numerical trace constant]\label{rem:constants}
One may take $C_{\mathrm{tr}}=143$ in
Theorem~\ref{T10:conj:law}.  Here are rational bounds for the
remaining scalar terms in \eqref{rev5:eq:final-constant}.
Direct differentiation of $u$ and $z=v-v_0$ gives
\[
 \mathfrak c(u)=\frac{\pi}{8\sqrt2}
                   +\frac{\pi(2-\sqrt2)}{48},\qquad
 \mathfrak c(z)=\frac{\pi}{12\sqrt2},\qquad
 C_*=\frac{\pi(3+\sqrt2)}{24}<\frac35.
\]
For example, the derivative jumps of $u$ have absolute values
$\pi/\sqrt2$ and zero, and
$\int_0^1|u'''|=(\pi/2)^2(2-\sqrt2)$.
The first-derivative jump of $z$ vanishes, while its
second-derivative jump and $\int_0^1|z'''|$ both equal
$\sqrt2(\pi/2)^2$.

For the strip term, set $r=11/7>\pi/2$ and use its positive
power series.  The ratio of successive terms from index five
onwards is at most $r/6$, so
\[
 b(1/2)\le\sum_{j=0}^4\frac{r^j}{j!(2j+1)}
       +\frac{r^5}{5!\,11}\frac1{1-r/6}
 =\frac{152829401}{80385480}<\frac{191}{100}.
\]
The elementary bounds $\pi>31/10$, $e^2>7$,
$\log2<7/10$, $\log(5/2)<1$ and $\zeta(3)<5/4$ give
\[
 \Delta<\frac73,\qquad
 B_1=\frac{3+\log2}{\pi^2}+\frac{2\Delta}{\pi}
       <\frac{5450}{2883}<\frac{19}{10},\qquad
 B_2<\frac{229700}{923521}<\frac14.
\]
The exponential comparisons follow from positive finite Taylor
sums, and the zeta bound from
$1+1/8+\int_2^\infty x^{-3}\,dx=5/4$.
Thus the scalar portion of \eqref{rev5:eq:final-constant} is less
than
\[
 4\frac{191}{100}+2+\frac{200}{961}
      +\frac35+\frac{19}{10}+\frac12
 =\frac{308676}{24025}<13.
\]
Lemma~\ref{pubpair:lem:numerical-gauge} now gives
$2C_V+\text{scalar portion}<143$.  The same estimates give $A_{\rm band}<144$, $C_0<138$ and $B_2<1/4$;
these upper bounds may replace the corresponding constants in the
band-local and gap estimates.
These estimates concern the
whole-fiber transport and are conservative; no optimal constant
is asserted.
\end{remark}

Finally $\int_{\mathbb R}d(\omega)^j\,d\omega=2^{-j}B(1/2,j)/\pi$
checks the quadratic and quartic coefficients, while
\eqref{rev5:eq:odd-bound} independently controls every odd Lipschitz test.



